% !PRINTED 477
\weberpage{477}
%<<CONT
or $f_1(\sqrt{-m})$. In this way one finds very easily, if one takes
$f(i)$, $f_1(\sqrt{-2})$, $f(\sqrt{-3})$, $f(\sqrt{-7})$ from the
formulae of the preceding section:
\weq{2}{f_1(\sqrt{-4})^8 = 8,}
\weq{3}{f_1(\sqrt{-16})^8 = 8\sqrt2(1+\sqrt2)^2,}
\weq{4}{f_1(\sqrt{-8})^8 = 8(1+\sqrt2),}
\weq{5}{f_1(\sqrt{-32})^8 = 32(1+\sqrt2)^2+8\sqrt2\sqrt{8(1+\sqrt2)^4+(1+\sqrt2)}.}
If we set:
\[
f_1(\sqrt{-32})^8 = 8x,
\]
then $x$ is a root of the biquadratic class equation:
\weq{6}{(x^2-24x-2)^2-8(8x+1)^2=0,}
which by the adjunction of $\sqrt2$ splits into two quadratic equations
\[
x^2-8(1\pm\sqrt2)^2x-2(1\pm\sqrt2)=0
\]
Furthermore we find in this way:
\weq{7}{f_1(\sqrt{-12})^4 = 2\sqrt[6]{2}(1+\sqrt3),}
\weq{8}{f_1(\sqrt{-28})^4 = 2\sqrt2(3+\sqrt7).}

We shall apply this method to other cases later on.

§ 130.

\gesperrt{Calculation of Class Invariants from the Schlaefli Modular Equations.}

I. If in one of the equations between $u=f(\omega)$
and $v=f(n\omega)$, or between $u=f_1(\omega)$, $v=f_1(n\omega)$ (§ 73),
we substitute for $u$ one of the known values of $f(\sqrt{-m})$ or $f_1(\sqrt{-m})$,
then we obtain an equation which $f(\sqrt{-n^2m})$
or $f_1(\sqrt{-n^2m})$ satisfies.

Under certain circumstances this equation still contains extraneous factors, which
have to be sought out and removed.

If $n$ is a prime, then by § 123 we obtain complete
information about these superfluous factors. If $n$ divides $m$,
then in the equation in question a linear factor belonging to the determinant $-m$
is to be separated off; if $-m$ is a quadratic non-residue
of $n$, then superfluous factors are not present at all;
if $-m$ is a quadratic residue of $n$, then a factor belonging to the discri-
%CONT>>
